% =========================================================================
% Chapter 1: Introduction
%
% Student Guide:
% - Introduces your topic, the problem, your objectives and how the thesis
%   is organized.
% - \section{...} makes a main heading and \subsection{...} a smaller one.
% - Give a heading a label, e.g. \label{sec:my-label}, and point to it from
%   the text with \ref{sec:my-label}.
% - Cite sources from 'references.bib' with \cite{key}.
% =========================================================================

\chapter{Introduction}
\label{ch:introduction}

The introduction chapter should start with the broad context of your research, giving background, an overview of the topic, and the motivation for the study \cite{leicester2023}. It should then narrow to the specific focus, clearly stating the problem, objectives, and research questions. After that, it should define the scope and limitations to set boundaries, highlight the significance of the research to show its value, and finally provide a roadmap of the thesis so readers understand its structure.

\section{Background}
\label{sec:background}

This section presents the broader context of your research area. Begin by explaining the general field of Computer Science and Engineering in which your study is situated, such as networking, cybersecurity, artificial intelligence, software engineering, etc. Provide historical developments, current trends, and key concepts that help the reader understand the environment in which your research problem exists. The background should establish a foundation that makes your research topic relevant and comprehensible.

\section{Overview}
\label{sec:overview}

Here, offer a concise overview of the specific topic you are addressing. Summarize the state of the art, highlight existing approaches, and explain how your research fits into the larger body of knowledge. This section bridges the gap between the general background and the specific problem, showing readers how your work connects to ongoing developments in the field.

\section{Motivation}
\label{sec:motivation}

Clearly articulate why the research is important. Discuss challenges in the field, unresolved issues, or technological gaps that motivate your study. Emphasize the potential impact of solving the problem, whether it advances theory, improves system performance, or addresses societal needs. This section should persuade readers that your research is both timely and necessary.

\section{Problem Statement}
\label{sec:problem-statement}

Define the specific problem your thesis addresses. Be precise and avoid vague language. Explain what is missing in current literature or practice and why this problem requires attention. The problem statement should be strong enough to justify the need for your research and serve as the anchor for your objectives and research questions.

\section{Objectives}
\label{sec:objectives}

List the main objectives of your study. These should be measurable and aligned with the problem statement. Distinguish between the overall aim (broad goal) and specific objectives (steps to achieve the aim). Objectives provide clarity on what your research intends to accomplish.

\section{Research Questions}
\label{sec:research-questions}

Formulate clear research questions that guide your investigation. These should directly stem from the problem statement and objectives, and they must be answerable within the scope of your thesis. Well-defined questions help structure your methodology and analysis.

\section{Scope and Limitations}
\label{sec:scope-limitations}

Define the boundaries of your research. Specify what is included, such as algorithms tested, datasets used, or environments considered, and acknowledge what is excluded. Discuss limitations such as resource constraints, assumptions, or methodological restrictions. This section ensures transparency and sets realistic expectations for your study.

\section{Significance of the Research}
\label{sec:significance}

Explain the contribution of your study. Highlight theoretical significance (advancing knowledge), practical significance (improving systems or applications), and potential societal or industrial impact. This section should convince readers of the value of your work and demonstrate how it contributes to the field of Computer Science and Engineering.

\section{Thesis Outline}
\label{sec:thesis-outline}

Conclude the introduction with a roadmap of the thesis. Briefly describe the content of each chapter, such as literature review, methodology, results, discussion, and conclusion. This helps readers navigate the document and understand its structure.

The following section shows how to use this template for preparing your thesis report.

\section{Headings}
\label{sec:headings}

In this \LaTeX{} template, section and subsection commands replace Word styles. Use the command \texttt{\textbackslash section\{...\}} for B-Head subsections and \texttt{\textbackslash subsection\{...\}} for C-Head sub-subsections.

Sub-subsection titles will automatically be added to the Table of Contents, directly below the previous sub-section title, each time you compile your document.

More guidelines are provided in the upcoming chapters of this template. Read them carefully. These guidelines are intended only to help you to use this thesis report template and should not be included in your final thesis report.
